\documentclass{aevidence}

% ==================================================
% NUR DIESE ANGABEN MUSST DU NORMALERWEISE ÄNDERN
% ==================================================

\journalvolume{14}
\journalissue{2}
\journalyear{2026}
\articledate{27.11.2026}
\articlepages{35--51}

\title{Titel des wissenschaftlichen Beitrags}

\subtitle{Optionaler Untertitel}

\author{
  Dr. Vorname Nachname
  \and
  Prof. Dr. Zweite Person
}

\shorttitlecommand{Kurztitel}

\articlekeywords{
  Begriff 1; Begriff 2; Begriff 3; Begriff 4
}

\abstracttext{
Hier steht die Zusammenfassung des Beitrags. Sie kann mehrere Sätze
umfassen und wird automatisch im Stil des jeweiligen Journals gesetzt.
}

\begin{document}

% ==================================================
% AB HIER NUR NOCH DEN ARTIKELTEXT SCHREIBEN
% ==================================================

\section{Einleitung}

Hier beginnt der Artikel.

Du kannst ganz normalen Fließtext schreiben. LaTeX kümmert sich um
Satz, Umbruch, Seitenzahlen und die typografische Gestaltung.

\section{Fragestellung}

Hier steht die Fragestellung.

\subsection{Eine Unterfrage}

Auch Unterabschnitte sind möglich.

\section{Methodik}

Hier wird die Methodik beschrieben.

\section{Ergebnisse}

Hier stehen die Ergebnisse.

\section{Diskussion}

Hier wird diskutiert, was aus den Ergebnissen folgt.

\section{Fazit}

Hier steht das Fazit.

% ==================================================
% OPTIONAL: ABBILDUNG
% ==================================================

% \begin{figure}[t]
%   \centering
%   \includegraphics[width=\linewidth]{figures/abbildung.png}
%   \caption{Eine bemerkenswerte Beobachtung.}
%   \label{fig:beobachtung}
% \end{figure}

% ==================================================
% OPTIONAL: LITERATURVERZEICHNIS
% ==================================================

% \begin{thebibliography}{9}
%
% \bibitem{muster}
% Muster, M. (2026).
% Ein Beitrag zur empirischen Plausibilität.
% \textit{Archiv für Ausreichende Evidenz}, 14(2), 12--19.
%
% \end{thebibliography}

\end{document}
